% ==========================================
% F4 - THU LOI EBGP TREN L4
% ==========================================

\section{F4 -- Thử lỗi BGP trên liên kết L4}
\label{sec:F4}

\subsection{Mục tiêu}

Thử nghiệm này kiểm tra khả năng chuyển
sang ISP C khi liên kết L4 giữa A2 và B1 bị lỗi.
Nhóm so sánh hai trường hợp:
chỉ tắt cổng phía A2 và tắt cổng ở cả hai đầu.
Thời gian được đo từ lúc gây lỗi
đến khi ping từ LAN-A tới LAN-C thông trở lại.

\subsection{Cách thực hiện}

Trước khi gây lỗi, mạng đã hội tụ
và chính sách ở Phần E đang ưu tiên ISP B.
Nhóm chạy ping liên tục từ A1,
dùng Loopback1 làm địa chỉ nguồn:

\begin{lstlisting}
ping 10.171.126.129 source Loopback1 repeat 1000 timeout 1
\end{lstlisting}

Đồng thời, nhóm kiểm tra đường tới LAN-C
trên A3 bằng lệnh:

\begin{lstlisting}
show ip bgp 10.171.126.128
\end{lstlisting}

\partzeroimage
{images/testing/F4/A3_bgp_before.png}
{Đường tới LAN-C trên A3 trước khi gây lỗi L4}
{fig:F4-bgp-before}

Hình \ref{fig:F4-bgp-before} ghi nhận
hai đường tới LAN-C.
A3 chọn đường qua A2 làm best path,
với next-hop \texttt{10.171.26.2}.
Tuyến này có LOCAL\_PREF 200
và AS\_PATH \texttt{64726 64826}.
Đường trực tiếp qua C1 \texttt{10.71.26.18}
có LOCAL\_PREF 100, đóng vai trò dự phòng.

\subsubsection{Chỉ tắt đầu A2}

Nhóm tắt FastEthernet1/0 trên A2,
là cổng nối sang B1 qua L4:

\begin{lstlisting}
! A2
conf t
interface FastEthernet1/0
 shutdown
end
\end{lstlisting}

Trong GNS3/Dynamips, tắt cổng phía A2
không làm cổng phía B1 tự chuyển sang down.
Nhóm ghi lại mốc gây lỗi và mốc ping
xuất hiện phản hồi trở lại trên đồng hồ của bản ghi.

\partzeroimage
{images/testing/F4/one_end_failure_start.png}
{Mốc gây lỗi khi chỉ tắt L4 phía A2: 00:00:52.972}
{fig:F4-one-start}

\partzeroimage
{images/testing/F4/one_end_ping_recovered.png}
{Mốc ping thông trở lại khi chỉ tắt đầu A2: 00:03:49.848}
{fig:F4-one-recovered}

Khoảng thời gian quan sát được là
\textbf{176,876 giây}.
Sau lần thử, cổng trên A2 được bật lại
bằng \texttt{no shutdown}; nhóm chờ
phiên eBGP được thiết lập và đường qua B
được ưu tiên trở lại trước khi thử tiếp.

\subsubsection{Tắt cả hai đầu L4}

Nhóm thực hiện lại thử nghiệm,
lần này tắt cả FastEthernet1/0 trên A2
và FastEthernet0/0 trên B1:

\begin{lstlisting}
! A2
conf t
interface FastEthernet1/0
 shutdown
end

! B1
conf t
interface FastEthernet0/0
 shutdown
end
\end{lstlisting}

\partzeroimage
{images/testing/F4/both_ends_failure_start.png}
{Mốc bắt đầu thử lỗi L4 ở cả hai đầu: 00:00:19.331}
{fig:F4-both-start}

\partzeroimage
{images/testing/F4/both_ends_ping_recovered.png}
{Mốc ping thông trở lại khi tắt cả hai đầu: 00:00:20.381}
{fig:F4-both-recovered}

Khoảng thời gian đo được là
\textbf{1,050 giây}.
Sau thử nghiệm, cả hai cổng được bật lại
bằng \texttt{no shutdown} và mạng được
đưa về trạng thái ổn định.

\subsection{Kết quả và giải thích}

\begin{table}[H]
\centering
\caption{Thời gian ping phục hồi trong thử nghiệm F4}
\label{tab:F4-results}
\small
\renewcommand{\arraystretch}{1.3}
\begin{tabularx}{\textwidth}{X c c c}
\toprule
\textbf{Trường hợp} &
\textbf{Mốc gây lỗi} &
\textbf{Ping trở lại} &
\textbf{Thời gian (s)} \\
\midrule
Chỉ tắt đầu A2 & 00:00:52.972 & 00:03:49.848 & 176,876 \\
Tắt cả hai đầu & 00:00:19.331 & 00:00:20.381 & 1,050 \\
\bottomrule
\end{tabularx}
\end{table}

Các mốc trong Bảng \ref{tab:F4-results}
được đọc từ đồng hồ hiển thị trong ảnh,
không phải timestamp hệ thống của router.
Thời gian phục hồi được tính bằng
mốc ping trở lại trừ mốc gây lỗi.
Các mốc thuộc những lần ghi riêng,
nên mỗi hàng được tính độc lập.
Đây là kết quả của từng lần đo,
không phải giá trị trung bình qua nhiều lần thử.

\partzeroimage
{images/testing/F4/A3_bgp_after.png}
{Đường tới LAN-C trên A3 sau khi mất đường qua ISP B}
{fig:F4-bgp-after}

Sau khi đường qua B bị rút,
Hình \ref{fig:F4-bgp-after} chỉ còn đường
qua C1 \texttt{10.71.26.18},
AS\_PATH \texttt{64826}, LOCAL\_PREF 100
và được chọn là best path.
Như vậy, AS A đã có đường dự phòng
đi ra LAN-C qua ISP C.

Khi chỉ tắt phía A2, A2 phát hiện
interface của mình bị down và có thể
ngắt phiên eBGP ngay nhờ cơ chế
BGP Fast External Fallover.
Tuy nhiên, cổng trên B1 vẫn có thể
ở trạng thái \texttt{up/up}.
B1 phải chờ cơ chế TCP/BGP phát hiện
mất kết nối, nên việc rút các route
không còn hợp lệ có thể chậm hơn.

Ping chỉ thành công khi cả đường đi
và đường về đều hoạt động.
Vì vậy, ngay cả khi A2/A3 đã có thể
đổi đường đi ra sang C, đường về LAN-A
vẫn có thể bị ảnh hưởng bởi route cũ
qua B1--A2 chưa được rút.
Đây là cách giải thích phù hợp với
thời gian gián đoạn gần Hold Time 180 giây;
hai ảnh bảng BGP không xác định riêng
thời điểm từng router đổi tuyến.
Do đó, 176,876 giây được ghi nhận là
thời gian ping phục hồi hai chiều,
không đồng nhất với thời gian riêng A2 chọn lại best path.

Khi tắt cả hai đầu, cả A2 và B1
đều thấy interface cục bộ bị down.
Fast External Fallover cho phép ngắt
phiên eBGP trực tiếp mà không cần chờ
Hold Time hết hạn, giúp các route
qua L4 được rút nhanh ở cả hai phía.
Điều này phù hợp với kết quả ping
phục hồi sau khoảng 1 giây.

\newpage
